\section{Conclusion}
\label{sec:conclusion}

Cancer-associated fibroblasts have emerged from the shadows of the tumour microenvironment to become a central focus of breast cancer research. The past decade---particularly the era of single-cell and spatial omics---has transformed our understanding of CAF biology, revealing a degree of heterogeneity and functional diversity that was unimaginable when these cells were first described as ``activated fibroblasts.''

We have organized this review around five interconnected themes that capture the current state of CAF research in breast cancer. Seven key messages emerge:

\begin{enumerate}
  \item \textbf{CAFs are heterogeneous, not homogeneous.} The myCAF/iCAF/apCAF classification provides a useful framework, but the true CAF landscape is far more complex, with multiple subpopulations and continuous transcriptional states shaped by local signalling gradients, ECM composition, and immune context.

  \item \textbf{CAF origins matter.} CAFs can arise from tissue-resident fibroblasts, bone marrow-derived mesenchymal stem cells, endothelial cells, and even tumour cells through EMT. These different origins may confer distinct functional properties and therapeutic vulnerabilities.

  \item \textbf{Spatial organization determines function.} The spatial distribution of CAF subpopulations---their positioning relative to tumour cells, immune cells, and vasculature---is a critical determinant of their biological effects. Spatial transcriptomics is revealing organized CAF niches that were invisible to bulk and dissociated single-cell analyses.

  \item \textbf{CAF--immune crosstalk is bidirectional and therapeutically actionable.} CAFs shape the immune landscape through paracrine signalling, physical barrier formation, and tolerogenic antigen presentation. Conversely, immune cells regulate CAF activation and heterogeneity. This crosstalk is a major determinant of immunotherapy efficacy.

  \item \textbf{CAFs are mediators of treatment resistance.} CAFs contribute to resistance to chemotherapy, endocrine therapy, immunotherapy, and targeted therapy through multiple mechanisms. Understanding these mechanisms is essential for developing combination strategies that overcome resistance.

  \item \textbf{CAF-targeted therapies are entering clinical trials.} Multiple strategies---including FAP-targeted antibodies, CAF reprogramming agents, and signalling pathway inhibitors---are being evaluated in clinical trials. The success of these approaches will depend on biomarker-driven patient selection and rational combination strategies.

  \item \textbf{Critical knowledge gaps remain.} The field lacks standardized CAF nomenclature, validated CAF-specific biomarkers, and prospective clinical trials demonstrating the efficacy of CAF-targeted therapies. Addressing these gaps requires community-wide coordination and investment.
\end{enumerate}

As spatial omics technologies mature and functional genomics approaches are applied to CAFs, we anticipate a new wave of discoveries that will further refine our understanding of CAF biology and open new therapeutic opportunities. The ultimate goal---translating CAF biology into improved outcomes for breast cancer patients---remains within reach, but will require the continued integration of basic, translational, and clinical research.
